%=====================================================================
\chapter{Web Payment Systems}\label{ch:payments}
%=====================================================================
\locator{4}{Part IV --- Money and Identity}

\begin{outcomes}
By the end of this chapter you will be able to:
\begin{tight}
  \item \textbf{Name} the parties to an online payment and \textbf{state} what
        each one does and is paid for.
  \item \textbf{Distinguish} the payment instruments available to a Pakistani
        shop, and \textbf{compare} them on cost, settlement time and risk.
  \item \textbf{Explain} who carries the loss when a payment goes wrong, for
        each instrument.
  \item \textbf{Compute} the true cost of cash on delivery, including the
        returns it causes.
  \item \textbf{Select} a payment mix for a stated shop, and \textbf{defend}
        the choice against the alternatives.
\end{tight}
\end{outcomes}

\begin{prereq}
\chref{checkout}: the order state machine, and that an order exists before it
is paid for. \chref{shipping}: order totals, so that percentages have something
to be percentages of.
\end{prereq}

\begin{keyterms}
payment instrument $\bullet$ card scheme $\bullet$ issuer $\bullet$ acquirer
$\bullet$ payment gateway $\bullet$ payment service provider $\bullet$ merchant
discount rate $\bullet$ settlement $\bullet$ chargeback $\bullet$ authorisation
$\bullet$ capture $\bullet$ mobile wallet $\bullet$ cash on delivery $\bullet$
return-to-origin $\bullet$ Raast $\bullet$ PayPak
\end{keyterms}

\section{What Happens When Somebody Pays}

A customer clicking \uiel{Pay} starts a conversation between five parties, four
of which the shop never sees. Knowing who they are is not trivia: it is how you
work out who to telephone when money is missing, and who will be out of pocket
if it stays missing.

\begin{definitionbox}[The five parties]
The \term{cardholder} or payer. Your customer.

\smallskip
The \term{issuer} --- the bank that gave them the card or holds their wallet.
It decides whether to approve the payment, and it is the party the customer
complains to.

\smallskip
The \term{scheme} --- Visa, Mastercard, PayPak --- which operates the rules and
the network between banks. It writes the rulebook that decides disputes.

\smallskip
The \term{acquirer} --- the bank that holds the \emph{shop's} account and
accepts the payment on its behalf. It carries the risk that the shop fails to
deliver, which is why it asks about your business before accepting you.

\smallskip
The \term{gateway} or \term{payment service provider} --- the software company
between the shop and the acquirer. It is the only one of the five you will
integrate with, and the only one whose documentation you will read.
\end{definitionbox}

\dsfig{%
\begin{tikzpicture}[
  pty/.style={rounded corners=3pt, draw=#1, line width=1.0pt, fill=#1!8,
              text=#1!60!black, font=\bfseries\scriptsize, align=center,
              text width=2.35cm, minimum height=0.95cm},
  arr/.style={-{Stealth[length=2.0mm]}, primary!55, line width=0.9pt},
  back/.style={-{Stealth[length=2.0mm]}, accent!60, line width=0.9pt,
               dashed},
  lbl/.style={font=\scriptsize, text=neutral, align=center, text width=2.9cm}
]
  \node[pty=secondary] (cu) at (0,1.6)    {Customer};
  \node[pty=greenhl]   (sh) at (3.0,1.6)  {Your shop};
  \node[pty=purplehl]  (gw) at (6.0,1.6)  {Gateway\\{\scriptsize\mdseries
       PSP}};
  \node[pty=neutral]   (ac) at (9.0,1.6)  {Acquirer};
  \node[pty=neutral]   (is) at (12.0,1.6) {Issuer};

  \foreach \a/\b in {cu/sh, sh/gw, gw/ac, ac/is}{\draw[arr] (\a) -- (\b);}
  \draw[back] (is.south) -- ++(0,-0.75) -| (sh.south)
       node[pos=0.25, below, font=\scriptsize, text=accent]
       {chargeback: the money comes back out};

  \node[lbl] at (1.5,2.75) {goods};
  \node[lbl] at (6.0,2.85) {\textbf{the only one you integrate with}};
  \node[lbl, text width=4.4cm] at (10.5,2.85)
       {\textbf{where the money actually is}, for the days before it settles};
  \node[font=\scriptsize\itshape, text=neutral, align=left, text width=12.0cm]
       at (5.2,-1.35)
       {You pay all four. The merchant discount rate is a single percentage
        that is then divided between scheme, acquirer and gateway --- which is
        why nobody will quote you its parts.};
\end{tikzpicture}}%
{The parties to a card payment, and the direction a chargeback travels. The
shop sits at one end of a chain it cannot see into, which is the central fact
about accepting cards.}{paymentparties}

\begin{conceptbox}[Authorisation is not payment]
Two separate events, and confusing them is how shops ship goods they are never
paid for.

\smallskip
\textbf{Authorisation} asks the issuer ``would this payment succeed?'' and
reserves the amount. It is instant, it is reversible, and it is not money.

\smallskip
\textbf{Capture} claims the reserved amount. Only now is the shop owed
anything.

\smallskip
\textbf{Settlement} is the acquirer actually paying the shop, typically days
later and net of fees.

\smallskip
A gateway that reports ``success'' may mean any of the three. Read its
documentation for which, because \chref{gateway}'s entire difficulty is that
your order status must mean the same thing as the gateway's.
\end{conceptbox}

\section{The Instruments a Pakistani Shop Can Actually Offer}

\begin{alertbox}[The platform ships no online payment method at all]
Checked rather than assumed. The payment extensions in the box are
\texttt{bank\_transfer}, \texttt{cheque}, \texttt{cod} and
\texttt{free\_checkout} --- \textbf{four methods, none of which takes money
online}. Every one of them records an intention and leaves the shop to collect
afterwards.

\smallskip
There is no card gateway, no wallet, nothing. Every real payment integration is
a third-party extension, bought or written, and \chref{gateway} is the chapter
that deals with that. A plan that assumed ``the platform handles payments'' has
already failed, and it is better to discover that here than in week fourteen.
\end{alertbox}

\rowcolors{2}{white}{lightbg}
\begin{longtable}{@{}L{2.9cm}L{2.5cm}L{2.4cm}L{3.0cm}L{3.1cm}@{}}
\toprule
\rowcolor{primary!15}
\textbf{Instrument} & \textbf{Typical cost} & \textbf{Settles} &
\textbf{Who loses if it goes wrong} & \textbf{Main difficulty} \\
\midrule
\endhead
Card (Visa, Mastercard) & 2--3\% & days & \textbf{the shop}, via chargeback &
disputes you cannot win \\
PayPak card & lower & days & the shop & domestic acceptance only \\
Mobile wallet (JazzCash, Easypaisa) & 1--2\% & days & mostly the shop &
reconciliation \\
Bank transfer / Raast & very low & near instant & effectively nobody &
confirming which order was paid \\
Cash on delivery & 1--3\% plus returns & on remittance & \textbf{the shop},
heavily & refusal at the door \\
\bottomrule
\end{longtable}

\smallskip
\textbf{Rates are indicative and negotiated.} Any figure here is a starting
point for a conversation with a provider, not a price. What does not change is
the \emph{shape}: online card acceptance is the most expensive and carries the
dispute risk; bank transfer is the cheapest and carries an operational burden
instead; cash on delivery looks cheap and is not.

\section{Cash on Delivery Is Not Free}

Cash on delivery dominates Pakistani e-commerce for a reason that has nothing
to do with technology: it asks the customer to trust nobody. They pay when the
goods are in their hand. A shop that refuses COD is refusing most of its
market.

It is also the most expensive instrument on the list, and the cost is invisible
because it does not arrive as a percentage.

\begin{worked}
A shop sells an item for \textbf{Rs~3{,}000} that costs it \textbf{Rs~2{,}100}
--- a 30\% gross margin, Rs~900 per order.

\smallskip
\textbf{The visible cost.} The courier charges a COD service fee, say
\textbf{2\%} of the collected amount: $0.02 \times 3{,}000 = Rs~60$.

\smallskip
\textbf{The invisible cost.} Some proportion of COD orders are refused at the
door --- the customer is out, has changed their mind, or ordered on an impulse
that has passed. Call it \textbf{20\%}, which is not a pessimistic figure for
unverified COD orders. A refused parcel is \term{return-to-origin}: the shop
pays delivery \emph{and} return, say \textbf{Rs~250} each way, and gets the
item back with nothing to show.

\smallskip
\textbf{Per 100 orders:}
\begin{tight}
  \item 80 delivered: $80 \times (900 - 60 - 250) = 80 \times 590 =
        Rs~47{,}200$
  \item 20 refused: $20 \times (-250 - 250) = -Rs~10{,}000$
  \item \textbf{Net: Rs~37{,}200}, against Rs~72{,}000 of gross margin on the
        same 100 orders if they were prepaid and delivery were free.
\end{tight}

\smallskip
\textbf{The effective cost of COD here is 48\% of gross margin} --- against
2--3\% for a card. The refusal rate, not the service fee, is almost all of it:
the 20 refused orders cost Rs~10{,}000 directly and also consumed the
Rs~20{,}000 of margin they would have produced.

\smallskip
\textbf{The lever that matters} is therefore the refusal rate, not the fee.
Dropping refusals from 20\% to 10\% is worth more than negotiating the courier
fee to zero --- which is why shops that take COD invest in confirming orders by
telephone before dispatch, and why an order confirmation step that feels like
friction is, here, the single most profitable thing in the pipeline.
\end{worked}

\begin{conceptbox}[What this says about the ``frictionless checkout'' advice]
Most writing about e-commerce checkout comes from markets where payment is
taken before dispatch. There, every removed step is pure gain, and the advice
is sound.

\smallskip
Where the shop is paid on delivery, an unverified order is a liability rather
than an asset, and a checkout optimised purely for completion produces more of
them. \textbf{Advice imported without its assumptions is how a shop loses money
by following best practice.} \chref{ux} made the same point about page weight;
this is the version that costs real money.
\end{conceptbox}

\section{The Local Rails}

\begin{examplebox}[What a Pakistani shop is actually plugged into]
\begin{tight}
  \item \textbf{State Bank of Pakistan} authorises and supervises the operators
        and providers. A gateway you integrate with should be an authorised
        one, and asking for that authorisation is a reasonable first question.
  \item \textbf{1LINK} operates the domestic switch connecting the banks, and
        \textbf{PayPak} is the domestic card scheme that runs over it. A PayPak
        card is cheaper to accept and works only in Pakistan.
  \item \textbf{Raast}, the State Bank's instant payment system, settles bank
        to bank in seconds at negligible cost. Its person-to-merchant use is
        the development most likely to change the arithmetic in this chapter,
        because it removes both the card fee and the COD refusal problem at
        once.
  \item \textbf{JazzCash} and \textbf{Easypaisa} are the mobile wallets, each
        with a merchant gateway, and between them they reach customers who have
        no card and no bank account --- which in this market is the majority.
\end{tight}

\smallskip
\textbf{Two absences worth knowing.} PayPal does not operate for merchants in
Pakistan, and Stripe does not support Pakistan-registered businesses. A great
deal of the integration tutorial material you will find online assumes one or
the other, and none of it applies. Students regularly lose a week to this.
\end{examplebox}

\begin{alertbox}[Dated facts, and how to treat them]
Everything in this section is a fact about a particular year. Rates are
negotiated, operators are authorised and de-authorised, and Raast's merchant
rollout is still changing what is possible.

\smallskip
\textbf{The durable part is the set of questions:} who is authorised, what does
it cost as a percentage and as a fixed fee, how long until the money is in the
shop's account, who decides a dispute, and what proportion of customers can
actually use it. Ask those five of any provider in any year and the answer will
be current. Memorising this year's rates will not help you in 2029; knowing
that settlement time and dispute liability are negotiable terms will.
\end{alertbox}

\begin{pitfall}
\begin{tight}
  \item \textbf{Treating an authorisation as a payment} and dispatching goods
        against it.
  \item \textbf{Quoting a client a card fee without the chargeback risk.} The
        2--3\% is the predictable part; the disputes are not.
  \item \textbf{Costing cash on delivery at the courier's service fee.} The
        refusal rate is most of the cost and appears in no invoice.
  \item \textbf{Offering only cards in a market where most customers have
        none.}
  \item \textbf{Following a payment integration tutorial written for PayPal or
        Stripe.} Neither is available to you here.
  \item \textbf{Assuming the platform can take payments out of the box.} Its
        four methods all collect later.
  \item \textbf{Letting the customer choose the currency without deciding who
        carries the exchange risk.}
\end{tight}
\end{pitfall}

\begin{lab}[Cost your own shop's payment mix]
\begin{lstlisting}[language=bash]
# --- 1. See what you actually have ----------------------------------
cd code
docker compose exec shop sh -c \
  'ls extension/opencart/admin/controller/payment/'
# Four files. Which of them takes money online? Write the answer down.

# --- 2. Enable bank transfer and read what it does ------------------
# Admin: Extensions > Extensions > Payments > Bank Transfer > Install.
# Place an order with it. What does the customer see, and what does
# the order status become?

# --- 3. Compare that with cash on delivery --------------------------
# Place one of each and compare the two orders' status histories.

# --- 4. Build the COD model for YOUR product ------------------------
# Take one product in your shop. Find its price. Assume a cost of 70%
# of price, Rs 250 each way for delivery, a 2% COD fee, and a 20%
# refusal rate. Compute net margin per 100 orders.

# --- 5. Find the break-even refusal rate ----------------------------
# At what refusal rate does COD earn the shop nothing at all?
# Solve it, do not guess.

# --- 6. Price the alternative ---------------------------------------
# Repeat step 4 for a 2.5% card fee, zero refusals, and the same
# delivery cost one way only. Which is better, and by how much?

# --- 7. Decide, and write the decision down -------------------------
# One paragraph: which instruments your shop will offer, in what
# order on the checkout page, and what you will measure to find out
# whether you were right.
\end{lstlisting}

\smallskip
\textbf{What to notice.} Step 5 surprises people. The break-even refusal rate
is lower than most shops' actual refusal rate, which means a shop can run on
COD for a year, grow its order count every month, and lose money on every
single one.
\end{lab}

\begin{checkpoint}
\begin{tightnum}
  \item A gateway reports ``success'' and your shop marks the order Processing
        and dispatches it. Describe two distinct ways the shop can end up
        having shipped goods it is never paid for, and say what you would check
        in the gateway's documentation to rule each one out.
  \item Using the worked model, compute the net margin per 100 orders if the
        refusal rate falls from 20\% to 10\%. State the improvement as a
        percentage, and compare it with what the shop would gain by
        negotiating the courier's COD fee from 2\% to 0\%.
  \item A client wants to accept international cards for customers abroad.
        List, in order, the four things you would establish before writing any
        code, and say which of them is most likely to stop the project.
\end{tightnum}
\end{checkpoint}

\begin{chaptersummary}
\begin{tight}
  \item Five parties: customer, issuer, scheme, acquirer, gateway. You
        integrate with one and pay all of them through a single merchant
        discount rate whose parts nobody will itemise.
  \item Authorisation, capture and settlement are three different events. A
        gateway's ``success'' means one of them, and your order status must
        mean the same thing.
  \item The platform ships \textbf{no online payment method} --- bank transfer,
        cheque, cash on delivery and free checkout all collect later. Every
        gateway is a third-party extension.
  \item Instruments differ less in price than in who carries the loss. Cards
        put the dispute risk on the shop; bank transfer puts an operational
        burden on it instead.
  \item Worked: cash on delivery at a 20\% refusal rate costs 48\% of gross
        margin, almost all of it from refusals rather than from the service
        fee. The refusal rate is the lever; the fee is not.
  \item Checkout advice written for prepaid markets does not transfer. Where
        payment is on delivery, an unverified order is a liability and
        confirming it is the most profitable step in the pipeline.
  \item Locally: SBP authorises, 1LINK switches, PayPak is the domestic
        scheme, Raast settles instantly, and the wallets reach customers with
        no card. PayPal and Stripe are not available to you.
  \item Rates and operators date; the five questions --- authorised, cost,
        settlement time, who decides disputes, who can use it --- do not.
\end{tight}
\end{chaptersummary}

\begin{reviewq}
  \item \textbf{(MCQ)} Which party carries the loss when a card payment is
        successfully disputed? (a)~the issuer; (b)~the scheme; (c)~the
        acquirer; (d)~the shop.
  \item \textbf{(MCQ)} How many of the payment extensions shipped with the
        platform take money online? (a)~none; (b)~one; (c)~two; (d)~all four.
  \item \textbf{(Short)} Distinguish authorisation, capture and settlement, and
        say which one a shop should require before dispatching goods.
  \item \textbf{(Short)} Explain why a 2\% COD service fee understates the cost
        of cash on delivery, using one number from the worked model.
  \item \textbf{(Short)} Give two reasons a Pakistani shop might prefer PayPak
        acceptance to Visa and Mastercard, and one reason it cannot rely on it
        alone.
  \item \textbf{(Applied)} A shop takes 600 orders a month at an average value
        of Rs~4{,}000 and a 28\% gross margin, all on cash on delivery with an
        18\% refusal rate and Rs~300 of delivery cost each way. Model the
        effect of introducing prepayment with a 2.5\% fee that 35\% of
        customers adopt, and state the monthly difference, showing every step.
  \item \textbf{(Applied)} You are asked to decide whether to add a mobile
        wallet to a shop that currently takes only COD. Set out what you would
        need from the provider, what you would change in the shop, what you
        would measure for how long, and what result would make you remove it
        again.
\end{reviewq}
